\chapter{Conclusion and Future Work}
\label{ch:conclusion}

\section{Summary of Contributions}
\label{sec:contrib}

This Final Year Project asked how multi-jurisdictional AI healthcare regulatory requirements can be systematically synthesized, scored, and interactively compared. The work answers that question through four integrated contributions: (i)~a thematic state-of-the-art connecting SaMD, privacy, validation, transparency, ethics, surveillance, and liability; (ii)~a curated twenty-jurisdiction dataset with an explicit seven-theme scoring ontology; (iii)~a Retrieval-Augmented Generation (RAG) module with hybrid retrieval, grounded answers, and citation policy (Chapter~\ref{ch:rag}); and (iv)~an interactive React/Streamlit dashboard---open-sourced at \repoLink{} ---that turns those indicators into explorable maps, comparisons, trends, clinical use-case views, and a RAG Assistant.

Relative to the SMART objectives in Chapter~\ref{ch:intro}, the project delivered measurable theme scores and device-count indicators, an achievable local dashboard without proprietary APIs, a defence-oriented RAG technical module inside the methodology chapter, and a time-bound full report aligned with the aivancity PFE Template and Writing Guidelines (mandatory structure, IEEE bibliography, generative-AI declaration). Scientifically, the contribution is an open comparative ontology plus interactive analytics with grounded retrieval; operationally, the contribution is a briefing-ready artifact for researchers, policymakers, and developers.

Empirically, the analysis confirms convergence toward risk-based device regulation and GDPR-inspired data protection, while highlighting persistent gaps in lifecycle surveillance, transparency mandates, and liability clarity---especially outside the most resourced agencies. The EU AI Act stands out as a binding inflection point for high-risk healthcare AI~\cite{euaiact2024}. The United States illustrates high innovation throughput under a more sectoral governance mix~\cite{fda2021}. Emerging economies often legislate privacy before AI-specific SaMD capacity fully materializes~\cite{who2021}.

\section{Limitations and Areas for Improvement}
\label{sec:limitations}

Limitations include subjectivity in expert scoring, incomplete cross-lingual coverage, heterogeneous device-count reporting, and the still-early stage of supervised NLP automation. Context constraints inherent to an independent research project---reliance on public data only and zero cloud budget---bounded the depth of primary-law review in some jurisdictions and deferred live monitoring.

Areas for improvement follow directly: dual annotation for score reliability; inclusion of non-English primary sources; sub-national EU/US profiles; formal usability testing with compliance officers; and archival PDF/A packaging in the submission pipeline. The dashboard remains a decision-support research tool, not a substitute for qualified legal counsel or notified-body advice.

No new empirical findings are introduced in this chapter; all quantitative claims are grounded in Chapter~\ref{ch:results}.

\section{Perspectives and Future Work}
\label{sec:future}

Concrete next steps include:
\begin{enumerate}
  \item building a labeled clause corpus and evaluating transformer classifiers with reported F1 against HITL gold labels~\cite{devlin2019,vaswani2017};
  \item upgrading the offline lexical RAG demo to a full dense FAISS/Chroma index with reported nDCG@5~\cite{lewis2020rag,karpukhin2020};
  \item adding longitudinal score versioning keyed to legislative events;
  \item expanding country coverage and sub-national U.S./EU profiles;
  \item integrating authenticated regulatory RSS/API feeds with change alerts;
  \item conducting formal usability studies with regulators and medtech compliance officers;
  \item exploring privacy-preserving analytics if de-identified hospital compliance logs become available under a DPIA;
  \item packaging automated regression tests for table/figure regeneration and continuous PDF/A export.
\end{enumerate}

These directions preserve the project's core ethic: automation may accelerate review, but human expertise remains responsible for compliance claims. Societal impact, if the roadmap succeeds, would be faster comparative learning for safer AI deployment---not automated legal determination.

\section{Answers to the Research Question}
\label{sec:answer}

Returning to the research question in Chapter~\ref{ch:intro}: multi-jurisdictional AI healthcare requirements \emph{can} be systematically synthesized by combining thematic literature organization with a stable scoring ontology; they \emph{can} be scored with transparent ordinal rubrics tied to primary instruments; and they \emph{can} be interactively compared through open dashboards that preserve drill-down to narratives. The approach does not eliminate legal complexity, but it reduces the cost of first-pass comparative intelligence and creates a substrate for future NLP automation under HITL control.

\section{Transferable Lessons}
\label{sec:lessons}

Three transferable lessons emerge for AI and data-science practice: (1) governance problems benefit from explicit ontologies as much as predictive problems benefit from feature schemas; (2) interactive analytics can be a first-class research contribution when paired with reproducible data; (3) responsible AI engineering includes knowing when automation must stop and expert review must begin.

A fourth practical lesson concerns reporting: PFE guideline compliance (structure, abstracts, bibliography mix, and page budgets) is itself part of professional communication quality, not an afterthought.

\section{Comparison Against Initial Objectives}
\label{sec:objcheck}

Table~\ref{tab:objcheck} summarizes objective completion at a high level.

\begin{table}[htbp]
\centering
\caption{SMART objectives versus delivered outcomes}
\label{tab:objcheck}
\begin{tabularx}{\textwidth}{@{}p{3.2cm} X c@{}}
\toprule
\textbf{Objective} & \textbf{Delivered outcome} & \textbf{Status} \\
\midrule
Specific synthesis across 20 jurisdictions & Dataset + SotA + profiles & Met \\
Measurable seven-theme scores & Ordinal scorecard + figures & Met \\
Achievable local dashboard & Streamlit app + GitHub repo & Met \\
Relevant under EU AI Act era & Comparative discussion + gap analysis & Met \\
Time-bound full report & This FPR draft for July 2026 & Met \\
\bottomrule
\end{tabularx}
\end{table}

Partial items (supervised NLP training; live monitoring) are explicitly roadmaped rather than over-claimed, which itself is part of methodological integrity.

\section{Implications for Practice}
\label{sec:implpractice}

For developers, theme-gap views suggest sequencing evidence packages by target-market weaknesses (e.g., transparency dossiers for EU high-risk pathways). For policymakers, capacity-building in post-market and liability themes may yield safer scaling than approval acceleration alone. For hospitals, the dashboard offers a briefing layer before procurement deep-dives. In all cases, users must treat scores as comparative research indicators.

\section{Closing Statement}

Governing AI in healthcare requires comparative intelligence that is rigorous enough for academic scrutiny and usable enough for professional briefing. By coupling a positioned literature synthesis with open, interactive analytics, this project provides a reproducible foundation for that intelligence---and a clear roadmap for turning curated knowledge into continuously maintained regulatory awareness.
